\section{5. Physical Description}\label{physical-description}

\subsection{5.1 Physical Architecture}\label{physical-architecture}

GSEG is located at one geographical site (\texttt{GSEG-\allowbreak{}REQ-\allowbreak{}ARC-\allowbreak{}001}). It
comprises GSN as a separate satellite-communications layer, C\&C, DPS
and EWS services, mission supervision and Flight Dynamics (MCC), the
archive (DAS) and shared GIS infrastructure.

The GIS platform provides logical separation between critical services
and archival and operator workloads. The preliminary architecture
assumes two separate equipment failure domains, each with the
resources required to accommodate the full critical-service load
following the loss of the other domain. This is a capacity allocation
(P), not demonstrated failover effectiveness. GSN RF equipment and MCC
workstations remain physically separated from the shared compute hosts.

\begin{figure}[H]
\centering
% Logical equipment groups and exchanges from DDF 5.5.1--5.5.2.
% Named source sets summarize independent logical inputs, not extra hardware.
\resizebox{0.98\linewidth}{!}{%
\begin{tikzpicture}[
  >=Stealth,
  gsegEquipment/.style={draw=HeadingBlue,fill=TableRowLightBlue,
    rounded corners=2pt,align=center,font=\small,text width=2.3cm,
    minimum height=1.4cm,inner sep=3pt},
  gsegSupport/.style={gsegEquipment,fill=white,text width=2.5cm},
  gsegExternal/.style={gsegEquipment,fill=white,text width=3.1cm,
    minimum height=1.1cm},
  gsegInfrastructure/.style={gsegEquipment,text width=11.6cm,
    minimum height=1.15cm},
  gsegMission/.style={draw=HeadingBlue,line width=0.85pt},
  gsegData/.style={gsegMission,-{Stealth[length=2mm]}},
  gsegOperations/.style={draw=HeadingBlue,dash pattern=on 5pt off 2pt,
    line width=0.7pt,-{Stealth[length=1.8mm]}},
  gsegArchive/.style={draw=HeadingBlue,dash pattern=on 2pt off 2pt,
    line width=0.7pt,-{Stealth[length=1.8mm]}},
  gsegServices/.style={draw=HeadingBlue!75,densely dotted,line width=0.65pt},
  gsegChannelLabel/.style={align=center,font=\scriptsize,text=HeadingBlue,
    inner sep=1pt},
  gsegNote/.style={align=left,font=\scriptsize,text=HeadingBlue,inner sep=1pt}
]
% One geographical site, containing exactly the seven GSEG subsystems.
\draw[draw=TableBorderBlue,rounded corners=5pt,line width=0.65pt]
  (-6.4,-3.0) rectangle (6.4,6.1);
\node[font=\small\bfseries,text=HeadingBlue] at (0,5.8)
  {GSEG --- one geographical site};

\node[gsegExternal] (sseg) at (-4.8,7.2) {Space Segment\\SSEG};
\node[gsegExternal] (recipient) at (4.8,7.2)
  {Authorized recipient\\interface};
\node[gsegEquipment,fill=white] (gsn) at (-4.8,4.0)
  {\textbf{GSN}\\Satellite RF gateway};
\node[gsegEquipment] (cc) at (-1.6,4.0)
  {\textbf{C\&C}\\LEO selection\\Commands};
\node[gsegEquipment] (dps) at (1.6,4.0)
  {\textbf{DPS}\\Event assessment};
\node[gsegEquipment] (ews) at (4.8,4.0)
  {\textbf{EWS}\\Warning delivery};

% Primary mission exchanges. There is no C&C -> DPS substitute path.
\draw[gsegMission,Stealth-Stealth] (sseg.south) -- (gsn.north);
\draw[gsegMission,Stealth-Stealth] (gsn.east) -- (cc.west);
\draw[gsegData] ($(gsn.north)+(0.5,0)$)
  -- (-4.3,5.05) -- (1.6,5.05) -- (dps.north);
\node[gsegChannelLabel,text width=6.0cm] at (-1.35,5.3)
  {Two-LEO products + preliminary trajectory};
\draw[gsegData] (dps.east) -- (ews.west);
\draw[gsegData] (ews.north) -- (recipient.south);

\node[gsegSupport] (mcc) at (-1.6,0.9)
  {\textbf{MCC}\\Flight dynamics\\Supervision};
\node[gsegSupport] (das) at (4.8,0.9)
  {\textbf{DAS}\\Mission archive};

% Precomputed plans/orbits and later warning revisions use separate routes.
\draw[gsegOperations] (mcc.north) -- (cc.south);
\node[gsegChannelLabel,align=right,text width=1.1cm] at (-2.35,2.45)
  {Orbits /\\plans};
\draw[gsegOperations] ($(mcc.north)+(0.65,0)$)
  -- (-0.95,2.55) -- (4.8,2.55) -- (ews.south);
\node[gsegChannelLabel,text width=4.0cm] at (2.0,2.8)
  {Post-warning revisions only};

% Independent source lists explicitly define the logical collection channels.
% No collection channel is shown as a new subsystem or as routing through MCC.
\node[gsegChannelLabel,text width=2.75cm] at (-4.85,2.0)
  {\textbf{Status from}\\GSN / C\&C / DPS / EWS / GIS};
\draw[gsegOperations] (-6.0,0.9) -- (mcc.west);
\node[gsegChannelLabel,text width=3.4cm] at (1.65,1.95)
  {\textbf{Records from}\\GSN / C\&C / DPS / EWS / MCC};
\draw[gsegArchive] (0.8,0.9) -- (das.west);

% Shared infrastructure is a provider layer, never a warning-processing stage.
\node[gsegInfrastructure] (gis) at (0,-1.9)
  {\textbf{GIS --- Ground Infrastructure Subsystem}\\
   Shared compute, networks, time reference, power, cooling and monitoring};
\draw[gsegServices] (gis.north) -- (0,-1.05);
\draw[gsegServices] (-5.8,-1.05) -- (5.8,-1.05);
\node[gsegChannelLabel,text width=11.6cm] at (0,-0.82)
  {GIS services to GSN, C\&C, DPS, EWS, MCC and DAS};

% Orbital-information providers are external; the entry point is MCC.
\node[gsegExternal] (orbital) at (-4.8,-4.0)
  {Orbital-information\\providers};
\draw[gsegOperations] (orbital.north) -- (-4.8,-3.2)
  -- (-6.9,-3.2) -- (-6.9,-0.4) -- (-1.6,-0.4) -- (mcc.south);

\node[gsegNote,anchor=north west,text width=7.4cm] at (-1.3,-3.5) {
  Solid: primary mission exchanges. Long dashes: status, plans and authorized control.\\
  Short dashes: archive records. Dotted: shared GIS services.\\
  Named source lists denote independent logical inputs, not extra equipment,
  physical buses or prerequisites for the first warning.};
\end{tikzpicture}%
}

\caption{Figure 5-1. GSEG equipment groups and logical relationships.
Solid lines: principal mission products; dashed: supervision and
archiving; dotted: GIS support. The diagram does not define wiring or RF
redundancy.}
\end{figure}

EWS is the sole path for delivering warnings to authorized recipients.
MCC may initiate later updates but sends them to EWS, not directly to
the recipient.

Required availability boundary: the 99.99\% assessment includes
GSN RF equipment and ground gateways, GIS, and GSEG-controlled
connections and warning-delivery interfaces. Independent failures of
SSEG spacecraft or recipient systems \emph{beyond the agreed interface
boundary} are not counted against GSEG availability. The single site
remains a common-cause failure exposure that is not removed by dividing
GIS into two domains.

\subsection{5.2 Product Tree}\label{product-tree}

The product tree defines seven GSEG subsystems and five subordinate
product items: two within DAS and three within DPS. The GIS product
identifier is \texttt{123456.02.07}.

\begin{figure}[H]
\centering
% Seven subsystem nodes, with exactly two DAS and three DPS product nodes.
\resizebox{0.98\linewidth}{!}{%
\begin{tikzpicture}[
  root/.style={draw=HeadingBlue,fill=TableRowLightBlue,rounded corners=2pt,
    align=center,font=\small\bfseries,text width=3.0cm,minimum height=1.0cm},
  item/.style={draw=TableBorderBlue,fill=white,rounded corners=2pt,
    align=center,font=\small,text width=4.1cm,minimum height=1.0cm,inner sep=3pt},
  product/.style={item,fill=TableRowLightBlue!35,text width=4.0cm},
  link/.style={draw=HeadingBlue,line width=0.65pt}
]
\node[root] (gseg) at (-5.7,0) {Ground Segment\\CI-02};
\node[item] (das) at (-0.5,0) {\textbf{DAS}\\Data Archiving System};
\node[item] (cc) at (-0.5,-1.4) {\textbf{C\&C}\\Command and Control};
\node[item] (mcc) at (-0.5,-2.8) {\textbf{MCC}\\Mission Control Center};
\node[item] (dps) at (-0.5,-4.2) {\textbf{DPS}\\Data Processing System};
\node[item] (gsn) at (-0.5,-5.6) {\textbf{GSN}\\Ground Station Network};
\node[item] (ews) at (-0.5,-7.0) {\textbf{EWS}\\Early Warning System};
\node[item,fill=TableRowLightBlue!55] (gis) at (-0.5,-8.4)
  {\textbf{GIS}\\Ground Infrastructure Subsystem};
\draw[link] (gseg.east) -- (-3.4,0) -- (-3.4,-8.4);
\foreach \name in {das,cc,mcc,dps,gsn,ews,gis} {
  \draw[link] (-3.4,0 |- \name.west) -- (\name.west);
}
\node[product] (raw) at (5.0,0.65) {Raw Data Archive};
\node[product] (processed) at (5.0,-0.65) {Processed Data Archive};
\draw[link] (das.east) -- (2.4,0);
\draw[link] (2.4,0.65) -- (2.4,-0.65);
\draw[link] (2.4,0.65) -- (raw.west);
\draw[link] (2.4,-0.65) -- (processed.west);
\node[product] (software) at (5.0,-2.9) {Data Processing Software};
\node[product] (detection) at (5.0,-4.2) {Event Detection Module};
\node[product] (trajectory) at (5.0,-5.5) {Trajectory Estimation Module};
\draw[link] (dps.east) -- (2.4,-4.2);
\draw[link] (2.4,-2.9) -- (2.4,-5.5);
\foreach \name in {software,detection,trajectory} {
  \draw[link] (2.4,0 |- \name.west) -- (\name.west);
}
\end{tikzpicture}%
}

\caption{Figure 5-2. GSEG product tree: seven subsystems and
five separate DAS and DPS product items.}
\end{figure}

The subordinate DAS elements are Raw Data Archive and
Processed Data Archive. The subordinate DPS elements are
Data Processing Software, Event Detection Module and
Trajectory Estimation Module. Their respective WBS identifiers
remain \texttt{123456.02.01.01}, \texttt{.01.02}, \texttt{.04.01},
\texttt{.04.02} and \texttt{.04.03}.

\subsection{5.3 Specification Tree}\label{specification-tree}

The GSEG TRS (\texttt{EWAS-\allowbreak{}HSB-\allowbreak{}GSEG-\allowbreak{}TRS-\allowbreak{}02}) is the parent specification
for seven planned subsystem TRSs. Subsystem TRSs have not yet been
issued. DDFs and DDJFs define and justify the design, but are not
branches of the \emph{requirements specification tree}; agreed
interfaces are controlled in ICDs.

\begin{figure}[H]
\centering
% Requirements hierarchy only; no DDF or DDJF branches.
\resizebox{0.98\linewidth}{!}{%
\begin{tikzpicture}[
  parent/.style={draw=HeadingBlue,fill=TableRowLightBlue,rounded corners=2pt,
    align=center,font=\small,text width=6.0cm,minimum height=0.95cm},
  child/.style={draw=TableBorderBlue,fill=white,rounded corners=2pt,
    align=center,font=\small,text width=4.3cm,minimum height=0.75cm},
  link/.style={draw=HeadingBlue,line width=0.65pt}
]
\node[parent] (system) at (-3.8,0) {EWAS System Requirements};
\node[parent] (gseg) at (-3.8,-1.6)
  {GSEG TRS\\\texttt{EWAS-HSB-GSEG-TRS-02}};
\draw[link] (system.south) -- (gseg.north);
\node[child] (das) at (4.4,-2.5) {DAS TRS \textit{(planned)}};
\node[child] (cc) at (4.4,-3.6) {C\&C TRS \textit{(planned)}};
\node[child] (mcc) at (4.4,-4.7) {MCC TRS \textit{(planned)}};
\node[child] (dps) at (4.4,-5.8) {DPS TRS \textit{(planned)}};
\node[child] (gsn) at (4.4,-6.9) {GSN TRS \textit{(planned)}};
\node[child] (ews) at (4.4,-8.0) {EWS TRS \textit{(planned)}};
\node[child] (gis) at (4.4,-9.1) {GIS TRS \textit{(planned)}};
\draw[link] (gseg.east) -- (0.5,-1.6) -- (0.5,-9.1);
\foreach \name in {das,cc,mcc,dps,gsn,ews,gis} {
  \draw[link] (0.5,0 |- \name.west) -- (\name.west);
}
\end{tikzpicture}%
}

\caption{Figure 5-3. GSEG requirements specification hierarchy.}
\end{figure}

\subsection{5.4 Physical Elements and Redundancy
Boundaries}\label{physical-elements-and-redundancy-boundaries}

\textbf{GSN --- communications layer.} Antennas, RF equipment and
link-handling equipment are physically separated from the managed
computing platform. The antenna configuration, radio redundancy and
demonstration of visibility and connectivity from one site remain within
the GSN DDF and system analyses.

\textbf{GIS --- shared infrastructure.} GIS is responsible for compute
hosts, networks and switches, common time, power, cooling and technical
monitoring. It hosts critical C\&C/DPS/EWS services in two equipment
domains, separated from DAS/MCC workloads. Each domain must provide the
full critical capacity after loss of the other. Host counts, state
replication, failover details and physical separation are to be designed
at lower levels. The site itself remains a common failure domain.

\textbf{MCC and DAS --- non-critical equipment.} MCC uses separate
operator workstations, while DAS archiving workloads are separated from
compute resources used by the critical warning chain. Both subsystems
use shared GIS services.

\subsection{5.5 Interfaces}\label{interfaces}

\subsubsection{5.5.1 External Interfaces}\label{external-interfaces}

\begingroup\GSEGTableStyle{3}
\begin{longtable}[]{|P{0.25}|P{0.19}|P{0.56}|}
\hline\rowcolor{TableHeaderBlue}
\begin{minipage}[b]{\linewidth}\raggedright
\color{white}\bfseries\sffamily
Boundary
\end{minipage} & \begin{minipage}[b]{\linewidth}\raggedright
\color{white}\bfseries\sffamily
GSEG owner
\end{minipage} & \begin{minipage}[b]{\linewidth}\raggedright
\color{white}\bfseries\sffamily
Exchanged products
\end{minipage} \\
\hline
\endhead
\hline
\endlastfoot
GSEG--SSEG (\texttt{IF-\allowbreak{}EXT-\allowbreak{}01}) & GSN & Commands, telemetry, GEO
reports, LEO packages and the onboard preliminary trajectory. \\ \hline
Authorized warning recipients & EWS & First warnings, revisions,
cancellations, status and available acknowledgements. \\ \hline
Orbital-data providers & MCC (operational ownership) & Predictions,
orbital and conjunction data; onward transmission to SSEG through C\&C
and GSN, taking \texttt{IF-\allowbreak{}EXT-\allowbreak{}04} into account. \\*
\end{longtable}
\endgroup

The physical ingress point for orbital data, the external connection
topology and protocol details require controlled ICDs. The ISL network
belongs to SSEG; its mere presence does not demonstrate communications
availability from a single GSEG site.

\subsubsection{5.5.2 Internal Interfaces}\label{internal-interfaces}

\begingroup\GSEGTableStyle{2}
\begin{longtable}[]{|P{0.31}|P{0.69}|}
\hline\rowcolor{TableHeaderBlue}
\begin{minipage}[b]{\linewidth}\raggedright
\color{white}\bfseries\sffamily
Relationship
\end{minipage} & \begin{minipage}[b]{\linewidth}\raggedright
\color{white}\bfseries\sffamily
Main products or services
\end{minipage} \\
\hline
\endhead
\hline
\endlastfoot
GSN ↔ C\&C & GEO reports, LEO requests, authenticated commands,
orbital-data forwarding packages, and command/link status. \\ \hline
GSN → DPS & LEO Observation Packages, onboard preliminary trajectory and
reception metadata. \\ \hline
DPS → EWS & Refined trajectory, impact-area assessment, uncertainty and
the event qualification outcome for warning; a warning is
dispatched only for a qualifying outcome. \\ \hline
GSN / C\&C / DPS / EWS / GIS → MCC & Spacecraft, command, service and
event status, classification and warning-decision outcome, and
infrastructure status. \\ \hline
MCC → C\&C & Current orbital data, conjunction assessments, approved
plans and operational directives. \\ \hline
MCC → EWS & Authorized requests for confirmation, revision or
cancellation. \\ \hline
GSN / C\&C / DPS / EWS / MCC → DAS & Mission products, event history,
\textbf{decisions to issue or withhold warnings}, warning versions and
audit data. \\ \hline
GIS → other GSEG subsystems & Compute, networking, time, power, cooling
and infrastructure management. \\*
\end{longtable}
\endgroup

These relationships are logical and do not define distinct cables or
protocols. Data schemas, security boundaries, data rates and error
handling are governed by the relevant ICDs and subsystem TRSs.

\textbf{Common time-base consistency:} the applicable SSEG specification
provides UTC timestamps in observation products
(\texttt{SSEG-\allowbreak{}REQ-\allowbreak{}FUN-\allowbreak{}006}) and spacecraft time synchronization
(\texttt{SSEG-\allowbreak{}REQ-\allowbreak{}ADCS-\allowbreak{}004}). GSEG retains the observation timestamp
received from SSEG and distinguishes it from the GSN product-reception
time, which marks the GSEG latency-budget boundary. GIS provides common
time services. Timestamp encoding on individual interfaces and allowable
correlation error remain subject to ICD agreement and the accuracy
budget; \textbf{the 50 µs accuracy specified for SSEG is not assigned to
GSEG}.
